% ============================================================
% Section 3 — Answer-Signal Starvation (draft v1, 2026-08-29)
% 占位: tie 账本表与成本相关面板数字 = [K8-pending]
% Fig~\ref{fig:mechanism} 占位图环境在文末, 生成器 fig_mechanism.py 待写
% ============================================================

\section{Answer-Signal Starvation in Multi-Hop Search RL}
\label{sec:starvation}

\subsection{Setting and notation}
\label{sec:starve-setting}

A rollout is a trajectory
$\tau=(s_1,a_1,o_1,\ldots,s_T,a_T)$ in which actions are protocol steps:
a plan, search queries, and a terminal answer (or an abstention); each
search returns retrieved paragraphs as the observation. During training
the environment is a closed candidate pool per question; gold supporting
documents $\mathcal{G}$ are privileged training annotations and play no
role at test time (\S\ref{sec:setup}). Outcome-supervised systems assign
each rollout a return of the form
\begin{equation}
R_{\mathrm{base}}(\tau) \;=\; R_{\mathrm{ans}}(\tau)
\;-\; \lambda_{s} N_{\mathrm{search}}(\tau) \;-\; P_{\mathrm{aux}}(\tau),
\label{eq:base-reward}
\end{equation}
a terminal correctness reward minus search, token, and format costs.

\subsection{Group-relative learning and answer-tied groups}
\label{sec:starve-ties}

Group-relative estimators draw $K$ rollouts of the same question and
compare them against each other. With a leave-one-out
baseline~\citep{kool2019buy,ahmadian2024rloo}, the advantage of rollout
$i$ is
\begin{equation}
A_i \;=\; R_i \;-\; \frac{1}{K-1}\sum_{j\neq i} R_j ,
\label{eq:rloo}
\end{equation}
so any reward component that is constant within the group cancels
exactly. We call a group \emph{answer-tied} when all $K$ rollouts carry
the same terminal correctness label---every rollout wrong or every
rollout right. In an answer-tied group, $R_{\mathrm{ans}}$ is such a
constant: it contributes nothing to the within-group ordering,
regardless of how differently the rollouts behaved. The group is not
gradient-free---cost terms still differ across rollouts---but it is
\emph{answer-blind}: terminal correctness supplies zero discrimination.

How often does this happen? We sample $K$ rollouts per question at the
training temperature on a stratified question set and count tied groups
(protocol in \S\ref{sec:mechanism}). Table~\ref{tab:tie-ledger} shows
the result: ties are not an edge case but the dominant regime, and they
grow with hop depth---at four hops, [84]\% of groups are answer-tied,
[76]\% uniformly wrong.\footnote{[K8-pending] preliminary values; final
numbers use $K{=}8$ and the SFT initialization matched to training.}
The deeper the question, the less the terminal reward can say.

\begin{table}[t]
\centering
\small
\begin{tabular}{lccc}
\toprule
 & 2-hop & 3-hop & 4-hop \\
\midrule
Answer-tied groups & [75]\% & [71]\% & [84]\% \\
\quad all wrong    & [~~]\% & [~~]\% & [76]\% \\
\quad all right    & [~~]\% & [~~]\% & [~~]\% \\
\bottomrule
\end{tabular}
\caption{The tie ledger: share of $K$-rollout groups in which terminal
correctness provides no within-group ordering, by hop depth.
[K8-pending placeholders.]}
\label{tab:tie-ledger}
\end{table}

\subsection{The residual ordering prefers cheap failure}
\label{sec:starve-cost}

What ranks the rollouts of an all-failure group under
Eq.~\ref{eq:base-reward}? Substituting into Eq.~\ref{eq:rloo} with
$R_{\mathrm{ans}}$ constant,
\begin{equation}
A_i \;=\; -\lambda_{s}\bigl(N_i - \bar{N}_{-i}\bigr)
\;-\;\bigl(P_i - \bar{P}_{-i}\bigr),
\label{eq:cost-ordering}
\end{equation}
where $\bar{N}_{-i}$, $\bar{P}_{-i}$ are leave-one-out means. The
surviving preference is purely economic: the rollout that searched
less, generated less, and gave up sooner outranks the one that kept
pursuing evidence---even if the latter recovered most of the required
chain and the former recovered none. On the questions where persistence
is the very skill to be learned, the effective training signal points
the other way. Empirically, in all-failure groups from our outcome-only
training logs, the advantage correlates [negatively] with search count
($\rho = [~~]$; Fig.~\ref{fig:mechanism}A).\footnote{Our estimator is
plain leave-one-out without group standard-deviation normalization;
pathologies specific to $z$-scored advantages, where a vanishing group
deviation can amplify an arbitrary residual signal without
bound~\citep{wang2026darkroom}, do not apply directly. The qualitative
point is shared: whatever signal remains in an all-failure group is the
signal the group trains on.}

\subsection{Coverage restores a meaningful ordering}
\label{sec:starve-fix}

Let $C_i \in [0,1]$ be the fraction of the gold evidence set that
rollout $i$ retrieves (formal definition in \S\ref{sec:method}). Adding
$\beta\,C$ to the return contributes
\begin{equation}
A^{\mathrm{cov}}_i \;=\; \beta\Bigl(C_i - \frac{1}{K-1}\sum_{j\neq i}
C_j\Bigr) \;=\; \beta\,\frac{K}{K-1}\bigl(C_i - \bar{C}\bigr)
\label{eq:cov-advantage}
\end{equation}
to the advantage: a common coverage offset cancels, and what remains is
precisely a within-group ordering by evidence progress. This term is
non-zero exactly when coverage varies inside the group---and it is in
the all-failure groups, where the terminal reward is silent, that this
variation carries otherwise-unused information. How often coverage
actually varies there, and how the repair reshapes training, is the
subject of \S\ref{sec:mechanism}.

% ---- Fig 1 占位(fig_mechanism.py 待写: A=全错组内 outcome-only 按成本
% 排序 vs +coverage 按证据进展排序的实例/相关面板; B=tie 账本柱状) ----
\begin{figure}[t]
\centering
\fbox{\parbox{0.95\columnwidth}{\centering\vspace{2em}
[Figure~1 placeholder: (A) ordering inside an all-failure group,
outcome-only vs.\ coverage-augmented; (B) answer-tied group share by
hop depth.]\vspace{2em}}}
\caption{Answer-signal starvation and its repair. [Placeholder.]}
\label{fig:mechanism}
\end{figure}
